% Môn: Toán
% Lớp: 10
% Chương: Chương I. Mệnh đề và tập hợp
% Bài: Bài 2. Tập hợp và các phép toán trên tập hợp · Xác định đường biên và miền nghiệm
% Số câu: 56
% App đọc file này trực tiếp — sửa rồi Commit trên GitHub, bấm Đồng bộ GitHub .tex

% Bài 2. Tập hợp và các phép toán trên tập hợp



% ===== Câu 52 =====
% ID: T10C1B2-34-TN
% Mức: VD
\dangbt{Xác định đường biên và miền nghiệm}
\begin{ex}
% ID: T10C1B2-34-TN
% Mức: VD
Cho tập $X=\{ x \in \mathbb{N} |~ (x^2-4)(x-1)=0\}$. Tính số phần tử của $X$.
\choice
{$1$}
{$3$}
{\True $2$}
{$4$}
\loigiai{
Ta có
  $(x^2-4)(x-1)=0\Leftrightarrow$ hoặc{&x^2-4=0
&x-1=0}
 $\Leftrightarrow$ hoặc{&x=2
&x=-2 $\notin\mathbb{N}$
&x=1.} 
 Do $x\in \mathbb{N} \Rightarrow x \in \{ 2;1\}$. Vậy $X$ có $2$ phần tử.
}
\end{ex}

% ===== Câu 53 =====
% ID: T10C1B2-35-TN
% Mức: VD
\begin{ex}
% ID: T10C1B2-35-TN
% Mức: VD
Viết tập hợp $A=\{ x \in \mathbb{N} |~ (2x+1)(x^2-5x+6)=0\}$ bằng cách liệt kê.
\choice
{$\{ -1;2;3 \}$}
{\True $\{ 2;3 \}$}
{$\left\{ -\dfrac{1}{2};2;3 \right\}$}
{$\{ -1;2 \}$}
\loigiai{
Ta có
  (2x+1)$(x^2-5x+6)$=0 $\Leftrightarrow$
 hoặc{&2x+1=0
&x^2-5x+6=0} $\Leftrightarrow$
 hoặc{&x=-$\dfrac{1}{2}$
&x=2
&x=3.} 
 Do $x \in \mathbb{N} \Rightarrow A=\left\{ 2;3 \right\}$.
}
\end{ex}

% ===== Câu 54 =====
% ID: T10C1B2-36-TN
% Mức: VDC
\begin{ex}
% ID: T10C1B2-36-TN
% Mức: VDC
Cho hai tập hợp $A=\left\{x\in \mathbb{N}|-5\le 2x-\dfrac{1}{3}\le 3\right\},\,B=\{x\in \mathbb{Z}|3x+10\gt 1\}$. Khẳng định nào sau đây là đúng?
\choice
{$\mathrm{C}_B A=\{2;3;4;\ldots\}$}
{$\mathrm{C}_B A=\{-3;-2;-1;2;3;4;\ldots\}$}
{$\mathrm{C}_B A=\{0;1;2;3;4;\ldots\}$}
{\True $\mathrm{C}_B A=\{-2;-1;2;3;4;\ldots\}$}
\loigiai{
Ta có $A=\left\{x\in \mathbb{N}|-5\le 2x-\dfrac{1}{3}\le 3\right\}$.
$-5\le 2x-\dfrac{1}{3}\le 3-5+\dfrac{1}{3}\le 2x\le 3+\dfrac{1}{3}-\dfrac{14}{3}\le 2x\le \dfrac{10}{3}-\dfrac{7}{3}\le x\le \dfrac{5}{3}$
Vì $x\in\mathbb{N}$ nên $x\in\{0; 1\}$. Vậy $A=\{0; 1\}$.

Ta có $B=\{x\in \mathbb{Z}|3x+10\gt 1\}$.
$3x\gt 1-103x\gt -9x\gt -3$
Vì $x\in\mathbb{Z}$ nên $x\in\{-2; -1; 0; 1; 2; 3; \ldots\}$. Vậy $B=\{-2; -1; 0; 1; 2; 3; \ldots\}$.

Ta cần tìm khẳng định đúng.
Xét tập hợp $B \setminus A = \{x \mid x \in B \text{ và } x \notin A\}$.
$B \setminus A = \{-2; -1; 0; 1; 2; 3; \ldots\} \setminus \{0; 1\} = \{-2; -1; 2; 3; \ldots\}$.
Khẳng định $D$ là $B \setminus A = \{-2; -1; 2; 3; \ldots\}$.
}
\end{ex}
